% !TEX program = pdflatex
% arara: pdflatex: { shell: false }
% arara: pdflatex: { shell: false }
% arara: latexmk: { options: ["-c"] }

\documentclass[a4paper, 11pt]{article}
\let\rnc\renewcommand
\let\nc\newcommand

% ──────────────────────────────────────────────
%  Encodage & langue
% ──────────────────────────────────────────────
\usepackage[T1]{fontenc}
\usepackage[utf8]{inputenc}
\usepackage[french]{babel}

% ──────────────────────────────────────────────
%  Polices
% ──────────────────────────────────────────────
\usepackage{cabin}
\usepackage[math]{iwona}
\usepackage{inconsolata}
\rnc{\familydefault}{\sfdefault}

% ──────────────────────────────────────────────
%  Mise en page
% ──────────────────────────────────────────────
\usepackage[top=1.5cm, bottom=2cm, left=1.5cm, right=1.5cm, headheight=14pt]{geometry}

\usepackage{fancyhdr, lastpage}
\pagestyle{fancy}
\fancyhf{}
\fancyfoot[L]{\small\textit{\CoursTitre}}
\fancyfoot[R]{\small\thepage/\pageref*{LastPage}}
\rnc{\headrulewidth}{0pt}
\rnc{\footrulewidth}{0.4pt}
\fancypagestyle{plain}{%
	\fancyhf{}
	\fancyfoot[L]{\small\itshape\CoursTitre}
	\fancyfoot[R]{\thepage/\pageref*{LastPage}}
	\rnc{\headrulewidth}{0pt}
	\rnc{\footrulewidth}{0.4pt}
}

\usepackage{tocloft}
\rnc{\cftsecleader}{\cftdotfill{\cftdotsep}}

% \usepackage{titlesec}
% \titleformat{\section}{\large\bfseries}{\thesection.~}{0em}{\underline}
% \titleformat{\subsection}{\normalsize\bfseries}{\thesubsection.~}{0em}{\underline}
% \titleformat{\subsubsection}{\normalsize}{\thesubsubsection.~}{0em}{\underline}

% ──────────────────────────────────────────────
%  Couleurs (sélection de 15)
% ──────────────────────────────────────────────
\usepackage[dvipsnames, table]{xcolor}

% Couleurs du template
\definecolor{coulPrimaire}{HTML}{2C3E50}   % bleu ardoise foncé
\definecolor{coulAccent}{HTML}{2980B9}     % bleu moyen
\definecolor{coulFond}{HTML}{F4F6F8}       % gris très clair
\definecolor{coulBordure}{HTML}{AEB6BF}    % gris moyen
\definecolor{coulAlerte}{HTML}{C0392B}     % rouge
\definecolor{coulVert}{HTML}{27AE60}       % vert
\definecolor{coulOrange}{HTML}{E67E22}     % orange

% Couleurs additionnelles : palette panachée
\definecolor{amethyst}{rgb}{0.6, 0.4, 0.8}           % violet
\definecolor{auburn}{rgb}{0.43, 0.21, 0.1}            % brun-rouge
\definecolor{caribbeangreen}{rgb}{0.0, 0.8, 0.6}      % vert caraïbe
\definecolor{goldenrod}{rgb}{0.85, 0.65, 0.13}        % jaune doré
\definecolor{mulberry}{rgb}{0.77, 0.29, 0.55}         % rose foncé
\definecolor{tealblue}{rgb}{0.21, 0.46, 0.53}         % bleu-vert
\definecolor{terracotta}{rgb}{0.89, 0.45, 0.36}       % terre cuite
\definecolor{asparagus}{rgb}{0.53, 0.66, 0.42}        % vert sauge

% ──────────────────────────────────────────────
%  Maths
% ──────────────────────────────────────────────
\usepackage{amsfonts, amsmath, amssymb, mathtools}
\usepackage{mathrsfs}   % \mathscr
\usepackage{yhmath}     % \wideparen etc.
\usepackage{xfp}        % calculs en virgule flottante
\usepackage{numprint}   % \numprint
\usepackage{pifont}     % \ding
\usepackage{relsize}    % \mathlarger etc.
\usepackage{eurosym}    % \euro

% ──────────────────────────────────────────────
%  TikZ & graphiques
% ──────────────────────────────────────────────
\usepackage{tikz}
\usetikzlibrary{arrows.meta, calc, patterns, decorations.markings,
	shapes.geometric, positioning, backgrounds}
\usepackage{pgfplots}
\pgfplotsset{compat=1.18}
\usepackage{tkz-tab}

% ──────────────────────────────────────────────
%  Figures & tableaux
% ──────────────────────────────────────────────
\usepackage{graphicx}
\usepackage{wrapfig}
\usepackage{float}
\usepackage{booktabs, array, tabularx}
\newcolumntype{R}[1]{>{\raggedleft\arraybackslash}b{#1}}
\newcolumntype{L}[1]{>{\raggedright\arraybackslash}b{#1}}
\newcolumntype{C}[1]{>{\centering\arraybackslash}b{#1}}

% ──────────────────────────────────────────────
%  Code
% ──────────────────────────────────────────────
\usepackage{listings}
\lstset{
	basicstyle=\ttfamily\small,
	breaklines=true,
	frame=none,
	numbers=left,
	numberstyle=\tiny\color{coulBordure},
	commentstyle=\color{coulBordure}\itshape,
	keywordstyle=\color{coulAccent}\bfseries,
	stringstyle=\color{coulAlerte},
	showstringspaces=false,
	xleftmargin=1cm,
}

% ──────────────────────────────────────────────
%  Divers
% ──────────────────────────────────────────────
\usepackage{etoolbox}
\usepackage{lipsum}
\usepackage{fontawesome5}
\usepackage{multicol}
\usepackage{soul}
\usepackage{contour, ulem}
\rnc{\ULdepth}{3pt}
\rnc{\ULthickness}{1pt}
\contourlength{1pt}
\rnc{\underline}[1]{\uline{\phantom{#1}}\llap{\contour{white}{#1}}}
\usepackage{parskip}
\setlength{\parskip}{0.5mm}
\usepackage{enumitem}
\setlist[itemize]{label=\textbullet, topsep=2pt, itemsep=1pt, parsep=0pt}
\setlist[enumerate]{topsep=2pt, itemsep=1pt, parsep=0pt}

% ──────────────────────────────────────────────
%  Liens (chargé en dernier)
% ──────────────────────────────────────────────
\usepackage[
	colorlinks=true,
	linkcolor=coulAccent,
	urlcolor=coulAccent,
	citecolor=coulAccent,
	pdftitle={Notes de cours},
	pdfauthor={},
]{hyperref}

% ──────────────────────────────────────────────
%  tcolorbox
% ──────────────────────────────────────────────
\usepackage{tcolorbox}
\tcbuselibrary{skins, breakable, theorems}

\tcbset{
	styleBase/.style={
			enhanced, breakable, arc=2pt, boxrule=0.5pt,
			left=6pt, right=6pt, top=4pt, bottom=4pt,
			fonttitle=\bfseries\small,
		},
	styleLateral/.style={
			enhanced, breakable, frame hidden, opacityback=0,
			left=8pt, right=6pt, top=3pt, bottom=3pt, arc=0pt,
		},
}

\newtcolorbox[auto counter, number within=section]{defin}[2][]{
	styleBase, colback=coulFond, colframe=coulAccent,
	colbacktitle=coulAccent, coltitle=white,
	title={\faBook\enspace Définition\ifstrempty{#2}{}{~:~#2}}, #1}

\newtcolorbox[auto counter, number within=section]{theo}[2][]{
	styleBase, colback=white, colframe=coulPrimaire,
	colbacktitle=coulPrimaire, coltitle=white,
	title={\faLightbulb\enspace Théorème\ifstrempty{#2}{}{~:~#2}}, #1}

\newtcolorbox[auto counter, number within=section]{prop}[2][]{
	styleBase, colback=white, colframe=coulBordure,
	colbacktitle=coulBordure, coltitle=coulPrimaire,
	title={\faCheckCircle\enspace Propriété\ifstrempty{#2}{}{~:~#2}}, #1}

\newtcolorbox[auto counter, number within=section]{methode}[2][]{
	styleBase, colback=white, colframe=coulAccent,
	colbacktitle=coulAccent!20!white, coltitle=coulAccent,
	title={\faCog\enspace Méthode\ifstrempty{#2}{}{~:~#2}}, #1}

\newtcolorbox[auto counter, number within=section]{ex}[2][]{
	styleLateral, borderline west={2.5pt}{0pt}{coulBordure},
	colback=white, colbacktitle=white, coltitle=coulPrimaire,
	title={\underline{Exemple\ifstrempty{#2}{}{~:~#2}}}, #1}

\newtcolorbox[auto counter, number within=section]{exercice}[2][]{
	styleBase, colback=white, colframe=coulOrange,
	colbacktitle=coulOrange!15!white, coltitle=coulOrange!70!black,
	title={\faPen\enspace Exercice\ifstrempty{#2}{}{~:~#2}}, #1}

\newtcolorbox{correction}[1][]{
	styleLateral, borderline west={2.5pt}{0pt}{coulVert}, colback=white,
	before upper={\textcolor{coulVert}{\faCheck}\enspace\textbf{Correction.}\enspace}, #1}

\newtcolorbox{rem}[1][]{
	styleLateral, borderline west={2.5pt}{0pt}{coulAccent}, colback=coulFond,
	before upper={\textbf{Remarque.}\enspace}, #1}

\newtcolorbox{attention}[1][]{
	styleBase, colback=white, colframe=coulAlerte,
	colbacktitle=coulAlerte, coltitle=white,
	title={\faExclamationTriangle\enspace Attention}, #1}

\newtcolorbox{demo}[2][]{
	styleLateral, borderline west={2.5pt}{0pt}{coulBordure}, colback=white,
	colbacktitle=white, coltitle=coulPrimaire,
	title={\underline{Démonstration}\ifstrempty{#2}{}{~:~#2}}, #1}

\newtcolorbox{preuve}[2][]{
	styleLateral, borderline west={2.5pt}{0pt}{coulBordure}, colback=white,
	colbacktitle=white, coltitle=coulPrimaire,
	title={\underline{Preuve}\ifstrempty{#2}{}{~:~#2}}, #1}

\newtcolorbox{vide}[1][]{ styleLateral, colback=white, #1}

% ──────────────────────────────────────────────
%  Macros générales
% ──────────────────────────────────────────────

% Raccourcis typographiques
\nc{\hligne}{\begin{center}\rule{0.5\textwidth}{1pt}\end{center}}
\nc{\cimg}[2]{\begin{center}\includegraphics[width=#1]{#2}\end{center}}
\nc{\wfig}[2]{\begin{wrapfigure}{r}{#1}\includegraphics[width=#1]{#2}\end{wrapfigure}}

% Raccourcis listes
\def\be{\begin{enumerate}}
\def\ee{\end{enumerate}}
\def\bi{\begin{itemize}}
\def\ei{\end{itemize}}

% Raccourcis medskip
\let\ms\medskip

% ──────────────────────────────────────────────
%  Macros mathématiques
% ──────────────────────────────────────────────

% Ensembles
\nc{\R}{\ensuremath{\mathbb{R}}}
\nc{\D}{\ensuremath{\mathbb{D}}}
\nc{\Q}{\ensuremath{\mathbb{Q}}}
\nc{\Z}{\ensuremath{\mathbb{Z}}}
\nc{\N}{\ensuremath{\mathbb{N}}}
\nc{\C}{\ensuremath{\mathbb{C}}}

% Courbes
\nc{\Cf}{\ensuremath{\mathcal{C}_f}}
\nc{\Cg}{\ensuremath{\mathcal{C}_g}}

% Vecteurs, normes, valeurs absolues
\rnc{\vec}[1]{\ensuremath{\overrightarrow{#1}}}
\nc{\norm}[1]{\ensuremath{\left\lVert#1\right\rVert}}
\nc{\abs}[1]{\ensuremath{\left\lvert#1\right\rvert}}

% Parenthèses, crochets, accolades automatiques
\nc{\pa}[1]{\ensuremath{\left(#1\right)}}
\nc{\bra}[1]{\ensuremath{\left[#1\right]}}
\nc{\brac}[1]{\ensuremath{\left\{#1\right\}}}

% Fractions entre parenthèses / crochets
\nc{\pfrac}[2]{\ensuremath{\pa{\frac{#1}{#2}}}}
\nc{\bfrac}[2]{\ensuremath{\bra{\frac{#1}{#2}}}}

% Matrices
\nc{\mat}[1]{\ensuremath{\begin{matrix}#1\end{matrix}}}
\nc{\pmat}[1]{\ensuremath{\pa{\mat{#1}}}}
\nc{\bmat}[1]{\ensuremath{\bra{\mat{#1}}}}
\nc{\smat}[1]{\ensuremath{\begin{smallmatrix}#1\end{smallmatrix}}}
\nc{\psmat}[1]{\ensuremath{\pa{\smat{#1}}}}
\nc{\determinant}[1]{\ensuremath{\left|\mat{#1}\right|}}

% Coordonnées
\nc{\coord}[2]{\ensuremath{\begin{pmatrix}#1\\#2\end{pmatrix}}}
\nc{\coordl}[2]{\ensuremath{\left(#1~;~#2\right)}}
\nc{\OIJ}{\ensuremath{\left(O;I,J\right)}}
\nc{\vOIJ}{\ensuremath{\left(O;\vec{i},\vec{j}\right)}}

% Flèches logiques
\nc{\eqv}{\Leftrightarrow}
\nc{\rarr}{\rightarrow}  \nc{\larr}{\leftarrow}
\nc{\Rarr}{\Rightarrow}  \nc{\Larr}{\Leftarrow}
\nc{\Lrarr}{\Leftrightarrow}
\rnc{\iff}{\ensuremath{\Lrarr}}

% Inégalités corrigées
\let\le\leqslant
\let\ge\geqslant

% Équation alignée raccourcie
\nc{\eqn}[1]{\begin{align*}#1\end{align*}}

% Intégrale display
\nc{\dint}{\ensuremath{\displaystyle\int}}

% Dérivées partielles
\nc{\pd}{\partial}
\nc{\pdx}[1]{\ensuremath{\frac{\pd #1}{\pd x}}}
\nc{\pdy}[1]{\ensuremath{\frac{\pd #1}{\pd y}}}
\nc{\pdt}[1]{\ensuremath{\frac{\pd #1}{\pd t}}}

% Couleur dans les maths
\nc{\mtc}[2]{\textcolor{#1}{\ensuremath{#2}}}
\nc{\cbox}[2]{\colorbox{#1}{\ensuremath{#2}}}
\let\tc\textcolor

% Numprint
\nc{\np}[1]{\ensuremath{\numprint{#1}}}

% Cartes (pifont)
\nc{\club}{\text{\ding{168}}}
\nc{\spade}{\text{\ding{171}}}
\rnc{\diamond}{\textcolor{red}{\text{\ding{169}}}}
\nc{\heart}{\textcolor{red}{\text{\ding{170}}}}

% ──────────────────────────────────────────────
%  Commande titre
% ──────────────────────────────────────────────
\nc{\Entete}{%
	\begin{center}
		\vspace*{-5mm}
		{\color{coulPrimaire}\rule{\linewidth}{1pt}}\\[0.1cm]
		{\small \hfill \CoursClasse}\\[0.2cm]
		{\LARGE\bfseries\color{coulPrimaire}%
			\textcolor{coulAccent}{\small\faBookOpen}\enspace\CoursTitre\enspace\textcolor{coulAccent}{\small\faBookOpen}}\\[0.15cm]
		{\color{coulPrimaire}\rule{\linewidth}{0.4pt}}
	\end{center}%
	\tableofcontents
	\medskip
}

% ──────────────────────────────────────────────
%  Variables
% ──────────────────────────────────────────────
\nc{\CoursClasse}{2nde}
\nc{\CoursTitre}{Probabilités : document de test}

% ══════════════════════════════════════════════
\begin{document}
\Entete

% ══════════════════════════════════════════════
\section{Palette des 15 couleurs disponibles}
% ══════════════════════════════════════════════

Voici toutes les couleurs définies dans ce template, avec leur nom et un aperçu visuel.

\begin{center}
	\begin{tikzpicture}
		\def\sw{3.5}  % largeur du rectangle couleur
		\def\sh{0.55} % hauteur de chaque ligne
		\def\gap{0.7} % espacement vertical

		% Colonne gauche
		\foreach \col/\nom/\i in {
				coulPrimaire/coulPrimaire/0,
				coulAccent/coulAccent/1,
				coulFond/coulFond/2,
				coulBordure/coulBordure/3,
				coulAlerte/coulAlerte/4,
				coulVert/coulVert/5,
				coulOrange/coulOrange/6,
				terracotta/terracotta/7
			}{
				\fill[\col, draw=coulBordure, line width=0.3pt]
				(0, {-\i*\gap}) rectangle (\sw, {-\i*\gap+\sh});
				\node[right, font=\small\ttfamily, coulPrimaire]
				at (\sw+0.15, {-\i*\gap+\sh/2}) {\nom};
			}

		% Colonne droite
		\foreach \col/\nom/\i in {
				amethyst/amethyst/0,
				auburn/auburn/1,
				caribbeangreen/caribbeangreen/2,
				asparagus/asparagus/3,
				goldenrod/goldenrod/4,
				mulberry/mulberry/5,
				tealblue/tealblue/6
			}{
				\fill[\col, draw=coulBordure, line width=0.3pt]
				(9, {-\i*\gap}) rectangle ({9+\sw}, {-\i*\gap+\sh});
				\node[right, font=\small\ttfamily, coulPrimaire]
				at ({9+\sw+0.15}, {-\i*\gap+\sh/2}) {\nom};
			}
	\end{tikzpicture}
\end{center}
%
\begin{rem}
	Les couleurs \texttt{coulFond} et \texttt{coulBordure} étant très claires,
	leur rectangle apparaît presque blanc : c'est normal, elles sont pensées
	pour des fonds de boîtes.
\end{rem}

% ══════════════════════════════════════════════
\section{Démonstration des environnements}
% ══════════════════════════════════════════════

\begin{defin}{Expérience aléatoire}
	Une \textbf{expérience aléatoire} est une expérience dont on ne peut pas
	prévoir le résultat. L'\textbf{univers} $\Omega$ contient toutes les
	issues possibles.
\end{defin}

\begin{theo}{Formule de la réunion}
	Pour tous événements $A$ et $B$ :
	\[
		P(A \cup B) = P(A) + P(B) - P(A \cap B).
	\]
\end{theo}

\begin{demo}{}
	On décompose $A\cup B = (A\setminus B)\cup(A\cap B)\cup(B\setminus A)$
	en parties disjointes, puis on applique l'additivité de $P$.
\end{demo}

\begin{prop}{Complémentaire}
	$P(\bar{A}) = 1 - P(A)$.
\end{prop}

\begin{preuve}{$P(\bar A) = 1-P(A)$}
	$A\cup\bar{A}=\Omega$ et $A\cap\bar{A}=\emptyset$, donc
	$P(A)+P(\bar{A})=1$.
\end{preuve}

\begin{methode}{Calculer une probabilité}
	\be
	\item Justifier l'équiprobabilité.
	\item Dénombrer $|\Omega|$ et $|A|$.
	\item Calculer $P(A)=\dfrac{|A|}{|\Omega|}$.
	\ee
\end{methode}

\begin{ex}{Dé équilibré}
	$P(\text{pair}) = \dfrac{3}{6} = \dfrac{1}{2}$.
\end{ex}

\begin{exercice}{Dé pipé}
	$P(\text{6})=2p$, les autres valent $p$. Trouver $P(\text{pair})$.
\end{exercice}

\begin{correction}
	$5p+2p=1 \Rarr p=\frac{1}{7}$, donc
	$P(\text{pair})=\frac{1}{7}+\frac{1}{7}+\frac{2}{7}=\frac{4}{7}$.
\end{correction}

\begin{attention}
	L'équiprobabilité doit être \textbf{justifiée} : un dé pipé n'est pas
	équiprobable.
\end{attention}

\begin{rem}
	On peut aussi utiliser la macro \verb|\eqv| pour $\eqv$, ou
	\verb|\Rarr| pour $\Rarr$.
\end{rem}

\newpage

% ══════════════════════════════════════════════
\section{Démonstration des macros mathématiques}
% ══════════════════════════════════════════════

\subsection{Ensembles}

\[
	x \in \R, \quad n \in \N, \quad z \in \C, \quad p \in \Z, \quad
	r \in \Q, \quad d \in \D.
\]

\subsection{Vecteurs, normes, valeurs absolues}

\begin{wrapfigure}{r}{0.28\linewidth}
	\centering
	\begin{tikzpicture}[scale=1.1]
		\draw[-{Stealth}, coulPrimaire] (0,0) -- (2,0)
		node[right, font=\small] {$x$};
		\draw[-{Stealth}, coulPrimaire] (0,0) -- (0,2)
		node[above, font=\small] {$y$};
		\draw[-{Stealth[length=6pt]}, coulAccent, very thick]
		(0,0) -- (1.5,1.2) node[above right, font=\small] {$\vec{u}$};
		\draw[dashed, coulBordure] (1.5,0) -- (1.5,1.2) -- (0,1.2);
		\node[below, font=\tiny, coulPrimaire] at (0.75,0) {$x$};
		\node[left, font=\tiny, coulPrimaire] at (0,0.6) {$y$};
	\end{tikzpicture}
	\caption{Vecteur $\vec{u}$}
\end{wrapfigure}

\lipsum[1-4]

\[
	\vec{AB} = \coord{x_B - x_A}{y_B - y_A}, \qquad
	\norm{\vec{u}} = \sqrt{x^2+y^2}, \qquad
	\abs{x-3} \le 2.
\]

\[
	\vOIJ, \quad \coordl{x}{y}, \quad \OIJ.
\]

\subsection{Parenthèses et fractions automatiques}

\[
	\pa{\frac{a+b}{c}}, \quad
	\bra{\frac{x}{2}}, \quad
	\pfrac{a}{b+c}, \quad
	\bfrac{1}{n}, \quad
	\brac{x \in \R \mid x > 0}.
\]

\subsection{Matrices et déterminants}

\[
	A = \pmat{a & b \\ c & d}, \qquad
	\det(A) = \determinant{a & b \\ c & d} = ad - bc.
\]

\subsection{Intégrale et dérivées partielles}

On a, pour tout $x\in\R$ :

\begin{enumerate}
	\item plop
	\item $\int_0^1 f(x)\,\mathrm{d}x = \frac{1}{3}$
	\item $\pdx{f} = \frac{\pd f}{\pd x}$
	\item $\pdt{u} = \frac{\pd u}{\pd t}$
\end{enumerate}

\subsection{Équation alignée avec \texttt{\textbackslash eqn}}

\eqn{
	(a+b)^2 &= a^2 + 2ab + b^2 \\
	(a-b)^2 &= a^2 - 2ab + b^2 \\
	(a+b)(a-b) &= a^2 - b^2
}

\subsection{Couleurs dans les maths}

\[
	\mtc{coulAccent}{P(A \cup B)} = \mtc{coulVert}{P(A)} +
	\mtc{coulVert}{P(B)} - \mtc{coulAlerte}{P(A \cap B)}.
\]

\[
	\cbox{coulFond}{x^2 + y^2 = r^2} \qquad
	\tc{mulberry}{f(x) = x^3 - 3x}.
\]

\subsection{Numprint et symboles cartes}

\[
	\np{12345.6789} \qquad \club \quad \heart \quad \diamond \quad \spade.
\]

\newpage

% ══════════════════════════════════════════════
\section{Figures TikZ et pgfplots}
% ══════════════════════════════════════════════

\begin{ex}{Cercle trigonométrique}
	\begin{wrapfigure}{r}{0.5\linewidth}
		\centering
		\begin{center}
			\begin{tikzpicture}[scale=2.0]
				\draw[-{Stealth}, coulPrimaire] (-1.3,0) -- (1.3,0)
				node[right, font=\small] {$x$};
				\draw[-{Stealth}, coulPrimaire] (0,-1.3) -- (0,1.3)
				node[above, font=\small] {$y$};
				\draw[thick, coulAccent] (0,0) circle (1);
				\draw[coulOrange, thick] (0,0) -- ({cos(50)},{sin(50)});
				\fill[coulOrange!15] (0,0) -- (0.28,0) arc(0:50:0.28) -- cycle;
				\draw[coulOrange] (0.28,0) arc(0:50:0.28);
				\node[coulOrange, font=\small] at (0.2,0.1) {$\theta$};
				\filldraw[coulAlerte] ({cos(50)},{sin(50)}) circle (0.03)
				node[above right, font=\small] {$M(\cos\theta,\sin\theta)$};
				\draw[dashed, coulBordure]
				({cos(50)},0) -- ({cos(50)},{sin(50)}) -- (0,{sin(50)});
			\end{tikzpicture}
		\end{center}
		\caption{Cercle trigonométrique}
	\end{wrapfigure}
	\lipsum
\end{ex}

\begin{ex}{Courbe de $f(x)=x^2-2x$ avec pgfplots}
	\begin{center}
		\begin{tikzpicture}
			\begin{axis}[
					width=10cm, height=5.5cm,
					axis lines=center,
					xlabel={$x$}, ylabel={$f(x)$},
					xmin=-1.5, xmax=3.5, ymin=-1.8, ymax=4,
					grid=major, grid style={dotted, coulBordure},
					tick label style={font=\small}, samples=100,
				]
				\addplot[coulAccent, thick, domain=-1.2:3.2] {x^2-2*x};
				\addplot[coulAlerte, only marks, mark=*] coordinates {(1,-1)};
				\node[coulAlerte, font=\small, above right]
				at (axis cs:1,-1) {min $(1,-1)$};
			\end{axis}
		\end{tikzpicture}
	\end{center}
\end{ex}

% ══════════════════════════════════════════════
\section{Tableaux de signes et de variations}
% ══════════════════════════════════════════════

\begin{ex}{Signe de $(x-1)(x+3)$}
	\begin{center}
		\begin{tikzpicture}
			\tkzTabInit[lgt=3, espcl=2.5]
			{$x$/1, $(x+3)$/1, $(x-1)$/1, $f(x)$/1}
			{$-\infty$, $-3$, $1$, $+\infty$}
			\tkzTabLine{,-,z,+,,+}
			\tkzTabLine{,-,,-,z,+}
			\tkzTabLine{,+,z,-,z,+}
		\end{tikzpicture}
	\end{center}
\end{ex}

\begin{ex}{Variations de $g(x)=x^2-2x+3$, $g'(x)=2x-2$}
	\begin{center}
		\begin{tikzpicture}
			\tkzTabInit[lgt=3, espcl=3.5]
			{$x$/1, $g'(x)$/1, $g(x)$/2}
			{$-\infty$, $1$, $+\infty$}
			\tkzTabLine{,-,z,+}
			\tkzTabVar{+/$+\infty$, -/$2$, +/$+\infty$}
		\end{tikzpicture}
	\end{center}
\end{ex}

% ══════════════════════════════════════════════
\section{Code Python}
% ══════════════════════════════════════════════

\begin{ex}{Simulation : lancer de pièce}
	\begin{lstlisting}[language=Python]
import random

n = 1000
pile = 0

for i in range(n):
    if random.randint(0, 1) == 1:
        pile = pile + 1

frequence = pile / n
print("Frequence :", frequence)
	\end{lstlisting}
\end{ex}

\newpage

% ══════════════════════════════════════════════
\section{Expériences aléatoires à deux épreuves}
% ══════════════════════════════════════════════

\begin{defin}{Épreuves indépendantes}
	Deux épreuves sont dites \textbf{indépendantes} si le résultat de l'une
	n'a aucune influence sur le résultat de l'autre.
\end{defin}

\begin{methode}{Calculer une probabilité associée à une expérience à deux épreuves}
	Léa dépose dans un panier 4 chemisiers indiscernables au toucher :
	1 blanc, 1 rouge et 2 verts. Dans un autre panier, elle y dépose
	2 jupes : 1 blanche et 1 noire. Elle tire successivement et au hasard
	un chemisier du premier panier et une jupe du deuxième panier.

	Ces deux épreuves sont indépendantes.

	\begin{enumerate}
		\item Représenter la situation à l'aide d'un arbre pondéré.
		\item Calculer la probabilité $P_1$ d'obtenir deux vêtements blancs.
		\item Calculer la probabilité $P_2$ de ne pas obtenir un chemisier vert
		      et d'obtenir une jupe noire.
	\end{enumerate}
\end{methode}

\begin{correction}
	\textbf{1)} La probabilité de tirer un chemisier vert est
	$P(V) = \dfrac{2}{4} = \dfrac{1}{2} = 0{,}5$.

	\medskip
	\textbf{Arbre pondéré :}

	\begin{center}
		\begin{tikzpicture}[
			level 1/.style={sibling distance=3.5cm, level distance=2cm},
			level 2/.style={sibling distance=1.4cm, level distance=2cm},
			every node/.style={font=\small},
			edge from parent/.style={draw, -{Stealth[length=4pt]}}
			]
			\node {}
			child { node {B}
					child { node {Bl} edge from parent node[left, fill=white, inner sep=2pt] {$\frac{1}{2}$} }
					child { node {N} edge from parent node[right, fill=white, inner sep=2pt] {$\frac{1}{2}$} }
					edge from parent node[left, fill=white, inner sep=2pt] {$\frac{1}{4}$}
				}
			child { node {R}
					child { node {Bl} edge from parent node[left, fill=white, inner sep=2pt] {$\frac{1}{2}$} }
					child { node {N} edge from parent node[right, fill=white, inner sep=2pt] {$\frac{1}{2}$} }
					edge from parent node[left, fill=white, inner sep=2pt] {$\frac{1}{4}$}
				}
			child { node {V}
					child { node {Bl} edge from parent node[left, fill=white, inner sep=2pt] {$\frac{1}{2}$} }
					child { node {N} edge from parent node[right, fill=white, inner sep=2pt] {$\frac{1}{2}$} }
					edge from parent node[right, fill=white, inner sep=2pt] {$\frac{1}{2}$}
				};
		\end{tikzpicture}
	\end{center}

	\textbf{2)} Sur un chemin de branches, les probabilités se
	\textbf{multiplient} :
	\[
		P_1 = P(B\,;\,Bl) = \frac{1}{4} \times \frac{1}{2} = \frac{1}{8}
		= 0{,}125.
	\]
	La probabilité d'obtenir deux vêtements blancs est égale à $12{,}5\,\%$.

	\ms

	\textbf{3)} L'événement correspond aux issues $(B\,;\,N)$ et $(R\,;\,N)$.
	Les probabilités de plusieurs feuilles s'\textbf{additionnent} :
	$$
		\begin{array}{rcccl}
			P_2 & = & P(B\,;\,N)                     & + & P(R\,;\,N)                     \\
			    & = & \frac{1}{4} \times \frac{1}{2} & + & \frac{1}{4} \times \frac{1}{2} \\
			    & = & 0{,}125                        & + & 0{,}125\qquad = 0{,}25         \\
		\end{array}
	$$
	La probabilité est égale à $25\,\%$.
\end{correction}

% ══════════════════════════════════════════════
\section{Épreuve de Bernoulli}
% ══════════════════════════════════════════════

\subsection{Définitions}

\begin{defin}{Épreuve de Bernoulli}
	Une \textbf{épreuve de Bernoulli} est une expérience aléatoire à deux
	issues que l'on peut nommer \textbf{succès} et \textbf{échec}.
\end{defin}

\begin{rem}
	Au succès on peut associer le nombre $1$ et à l'échec le nombre $0$.
\end{rem}

\begin{ex}{Exemples d'épreuves de Bernoulli}
	\begin{enumerate}
		\item \textbf{Pile ou face :} on considère comme succès « obtenir pile »
		      et comme échec « obtenir face ».
		\item \textbf{Lancer d'un dé :} on considère comme succès « obtenir un 1 »
		      et comme échec « ne pas obtenir un 1 ». La loi de Bernoulli est :
		      \begin{center}
			      \begin{tabular}{|c|c|c|} \hline
				      $x_i$                             & $1$           & $0$           \\ \hline
				      \rule[-2mm]{0mm}{7mm}$P(X = x_i)$ & $\frac{1}{6}$ & $\frac{5}{6}$ \\ \hline
			      \end{tabular}
		      \end{center}
	\end{enumerate}
\end{ex}

\begin{defin}{Loi de Bernoulli}
	Une \textbf{loi de Bernoulli} de paramètre $p$ est la loi de probabilité
	d'une épreuve de Bernoulli vérifiant :
	\bi
	\item la probabilité d'obtenir $1$ est égale à $p$ ;
	\item la probabilité d'obtenir $0$ est égale à $1-p$.
	\ei

	\begin{center}
		\begin{tabular}{|c|c|c|} \hline
			$x_i$                             & $1$ & $0$     \\ \hline
			\rule[-2mm]{0mm}{7mm}$P(X = x_i)$ & $p$ & $1 - p$ \\ \hline
		\end{tabular}
	\end{center}
\end{defin}

\begin{ex}{Paramètres}
	Dans les exemples précédents :
	\begin{enumerate}
		\item $p = \dfrac{1}{2}$ \quad (pile ou face)
		\item $p = \dfrac{1}{6}$ \quad (obtenir un 1 avec un dé)
	\end{enumerate}
\end{ex}

\subsection{Espérance}

\begin{prop}{Espérance d'une loi de Bernoulli}
	Soit $X$ une variable aléatoire qui suit la loi de Bernoulli de
	paramètre $p$. Alors :
	\[
		E(X) = p.
	\]
\end{prop}

\begin{demo}{}
	\[
		E(X) = 1 \times p + 0 \times (1-p) = p.
	\]
\end{demo}

\begin{methode}{Reconnaître une situation modélisée par une loi de Bernoulli}
	Après la correction d'un contrôle, le professeur compte que 24 élèves
	ont obtenu une note supérieure ou égale à 10, et 6 ne l'ont pas obtenue.
	Le professeur choisit une copie au hasard. Justifier que cette situation
	peut être modélisée par une loi de Bernoulli.
\end{methode}

\begin{correction}
	Cette expérience aléatoire possède deux issues. La probabilité de
	l'issue « la copie indique une note $\ge 10$ » est :
	\[
		p = \frac{24}{30} = \frac{4}{5}.
	\]
	On considère cet événement comme le succès. La situation est donc
	modélisée par une loi de Bernoulli de paramètre $p = \dfrac{4}{5}$.
\end{correction}

\subsection{Répétitions d'épreuves de Bernoulli}

\begin{defin}{Expériences identiques et indépendantes}
	Plusieurs expériences sont \textbf{identiques et indépendantes} si :
	\bi
	\item elles ont les mêmes issues ;
	\item chaque issue possède la même probabilité.
	\ei
\end{defin}

\begin{ex}{Lancer de dé répété}
	On lance 5 fois de suite un dé à six faces. On considère comme succès
	« obtenir un six » ($p = \frac{1}{6}$) et comme échec « ne pas obtenir
	un six » ($1-p = \frac{5}{6}$).

	On répète ainsi 5 fois la même épreuve de Bernoulli de paramètre
	$p = \dfrac{1}{6}$, de manière indépendante.

	\begin{center}
		\begin{tikzpicture}
			\draw[-{Stealth}, coulPrimaire] (0,0) -- (0,1.2)
			node[above, font=\small] {Succès \quad $\frac{1}{6}$};
			\draw[-{Stealth}, coulPrimaire] (0,0) -- (0,-1.2)
			node[below, font=\small] {Échec \quad $\frac{5}{6}$};
			\filldraw[coulAccent] (0,0) circle (0.06);
		\end{tikzpicture}
	\end{center}
\end{ex}

\begin{ex}{Lancer de pièce répété}
	On lance 20 fois de suite une pièce. On considère comme succès
	« obtenir Pile » et comme échec « obtenir Face ».

	Ces épreuves de Bernoulli sont identiques et indépendantes avec
	paramètre $p = 0{,}5$.
\end{ex}

\begin{methode}{Calculer une probabilité associée à une épreuve de Bernoulli}
	Une urne contient 3 boules blanches et 2 boules rouges. On tire au
	hasard une boule et on la remet dans l'urne. On répète l'expérience
	deux fois de suite.

	\begin{enumerate}
		\item Représenter l'ensemble des issues dans un arbre.
		\item Déterminer les probabilités :
		      \begin{enumerate}
			      \item On tire deux boules blanches.
			      \item On tire une boule blanche et une boule rouge.
			      \item On tire au moins une boule blanche.
		      \end{enumerate}
	\end{enumerate}
\end{methode}

\begin{correction}
	On note $A$ : « tirer une boule blanche » et $\bar{A}$ : « tirer une
	boule rouge ».
	\[
		P(A) = \frac{3}{5} = 0{,}6 \qquad P(\bar{A}) = \frac{2}{5} = 0{,}4.
	\]

	\textbf{1) Arbre pondéré :}

	\begin{center}
		\begin{tikzpicture}[
			level 1/.style={sibling distance=4cm, level distance=2cm},
			level 2/.style={sibling distance=2cm, level distance=2cm},
			every node/.style={font=\small},
			edge from parent/.style={draw, -{Stealth[length=4pt]}}
			]
			\node {}
			child { node {$A$}
					child { node {$A$}
							node[below, coulAccent] {$0{,}36$}
							edge from parent node[left] {$0{,}6$} }
					child { node {$\bar{A}$}
							node[below, coulOrange] {$0{,}24$}
							edge from parent node[right] {$0{,}4$} }
					edge from parent node[left] {$0{,}6$}
				}
			child { node {$\bar{A}$}
					child { node {$A$}
							node[below, coulOrange] {$0{,}24$}
							edge from parent node[left] {$0{,}6$} }
					child { node {$\bar{A}$}
							node[below, coulAlerte] {$0{,}16$}
							edge from parent node[right] {$0{,}4$} }
					edge from parent node[right] {$0{,}4$}
				};
		\end{tikzpicture}
	\end{center}

	\textbf{2a)} Deux boules blanches $\Rarr$ issue $(A\,;A)$ :
	\[
		P_1 = 0{,}6 \times 0{,}6 = 0{,}36.
	\]

	\textbf{2b)} Une blanche et une rouge $\Rarr$ issues
	$(A\,;\bar{A})$ et $(\bar{A}\,;A)$ :
	\[
		P_2 = 0{,}24 + 0{,}24 = 0{,}48.
	\]

	\textbf{2c)} Au moins une blanche $\Rarr$ issues
	$(A\,;A)$, $(A\,;\bar{A})$ et $(\bar{A}\,;A)$ :
	\[
		P_3 = 0{,}36 + 0{,}24 + 0{,}24 = 0{,}84.
	\]

	\begin{rem}
		On peut aussi utiliser le complémentaire :
		$P_3 = 1 - P(\bar{A}\,;\bar{A}) = 1 - 0{,}16 = 0{,}84$.
	\end{rem}
\end{correction}

\end{document}
